\chapter*{Preface}
The series asks whether pure order and pure chance can be told apart, and why the primes look random yet obey laws. This floor is the section through the building. The drawing's title is the meeting: structure on one side, randomness on the other, and the members that only look like one of them. There is no next floor. What remains is labelled, and two projects stop on a line that is still open.

The first four parts ask the question in the form the earlier floors have made precise. A deterministic set can match a mean and a variance and still fail a short-interval test. That failure is Maier's theorem. A function built from the primes, $\log|\zeta(\tfrac12+it)|$, is nevertheless Gaussian at a random height. That is Selberg's theorem. Between them the book files the other meetings: Kubilius's model, a random multiplicative function, a sieve read as an expectation, a class group and a Selmer group modelled by a random group, and the form factor shared by a zero and an eigenvalue.

Where an earlier volume already contains the argument, it is cited and not repeated: the prime number theorem, Bombieri--Vinogradov, the sine kernel, Erd\H os--Kac, and Roth. A conjecture is labelled as a conjecture. Gallagher's Poisson law stays conditional on the uniform Hardy--Littlewood conjecture, and the hypothesis is in the heading.

The last part is the Research Workshop. It is the same meeting practised as a skill: how to read a paper about it, on Maier's theorem or on Selberg's, and two projects whose first steps are theorems of this volume and whose next steps are named and open.
